%---------------------------------------------------------------------
% Part five: Islamic history. Three questions.
%
% History answers are won on names and dates, so every date here is
% given in both Hijri and Gregorian and was checked individually. Where
% the traditional account is disputed by historians (Ibn al-Alqami's
% alleged treachery, for one) that is said plainly rather than passed on
% as settled fact.
%---------------------------------------------------------------------

\section*{حصہ پنجم --- اسلامی تاریخ}
\markboth{اسلامی تاریخ}{}

%---------------------------------------------------------------------
\begin{sawal}{سوال \textenglish{13} \ \textnormal{\small(خزاں \textenglish{2021}، سوال \textenglish{7} اور بہار \textenglish{2022}، سوال \textenglish{8})}}
اموی عہد کی معاشی و معاشرتی سرگرمیوں کا جائزہ پیش کریں، نیز حضرت عمر بن
عبدالعزیزؒ کی اصلاحات پر نوٹ لکھیں۔
\end{sawal}

\begin{hal}
\textbf{بنیادی معلومات۔} خلافتِ امویہ: \textenglish{41--132}ھ
(\textenglish{661--750}ء)۔ بانی: حضرت امیر معاویہؓ۔ دارالحکومت: دمشق۔ کل چودہ
خلفاء۔

\medskip
\textbf{(الف) معاشی سرگرمیاں:}
\begin{itemize}\setlength{\itemsep}{1pt}
  \item \textbf{زراعت} --- نہریں، بند، نظامِ آبپاشی؛ بنجر زمین آباد کرنے والے
        کو مالک بنانے کی پالیسی۔
  \item \textbf{خالص اسلامی سکہ} --- عبد الملک بن مروان نے \textenglish{77}ھ
        میں پہلا اسلامی دینار و درہم جاری کیا (پہلے رومی و ایرانی سکے چلتے
        تھے) --- اموی عہد کا سب سے بڑا معاشی کارنامہ۔
  \item \textbf{مالیاتی نظام} --- دیوانِ خراج کی تنظیم، عشر، خراج، جزیہ کی
        الگ الگ مدیں۔
  \item \textbf{ڈاک کا نظام (البرید)} --- حضرت امیر معاویہؓ نے قائم کیا، جس نے
        وسیع سلطنت کا انتظام ممکن بنایا۔
\end{itemize}

\medskip
\textbf{(ب) معاشرتی سرگرمیاں:}
\begin{itemize}\setlength{\itemsep}{1pt}
  \item \textbf{عربی کو سرکاری زبان بنانا} --- عبد الملک بن مروان کا اقدام۔
  \item \textbf{تعمیرات} --- قبۃ الصخرہ (بیت المقدس، \textenglish{72}ھ)، جامع
        مسجد دمشق، مسجدِ نبوی و حرام کی توسیع۔
  \item \textbf{صحت} --- ولید بن عبد الملک نے پہلا باقاعدہ دار الشفاء قائم
        کیا۔
  \item \textbf{کمزور پہلو} --- عرب و عجم کا امتیاز، موالی (غیر عرب مسلمانوں)
        سے دوسرے درجے کا سلوک --- یہی آگے چل کر عباسی تحریک کی بنیاد بنا۔
\end{itemize}

\medskip
\textbf{(ج) حضرت عمر بن عبدالعزیزؒ کی اصلاحات۔} دورِ خلافت:
\textenglish{99--101}ھ (\textenglish{717--720}ء) --- صرف ڈھائی سال، مگر انہیں
\textbf{''خلیفۂ راشدِ پنجم''} کہا جاتا ہے۔

\begin{enumerate}\setlength{\itemsep}{1pt}
  \item \textbf{مظالم کا خاتمہ} --- ناجائز مصادرہ کی گئی جائیدادیں مالکوں کو
        واپس کیں، ابتدا اپنے خاندان سے کی۔
  \item \textbf{منبروں کی اصلاح} --- خطبوں میں حضرت علیؓ پر سبّ کا سلسلہ بند
        کیا: \ar{إِنَّ اللهَ يَأْمُرُ بِالْعَدْلِ وَالْإِحْسَانِ} (النحل،
        \textenglish{90}) پڑھنے کا حکم دیا --- یہ روایت آج تک قائم ہے۔
  \item \textbf{تدوینِ حدیث کا سرکاری حکم} --- گورنروں کو لکھا کہ احادیث جمع
        کی جائیں --- حدیث کی تاریخ کا فیصلہ کن موڑ۔
  \item \textbf{ٹیکس کی اصلاح} --- نئے مسلمان ہونے والوں سے جزیہ ہٹا دیا۔
  \item \textbf{ذاتی سادگی} --- اپنی جائیداد بیت المال کو دے دی۔
\end{enumerate}

\smallskip
مشہور روایت ہے کہ ان کے عہد میں خوش حالی اس حد تک پہنچی کہ زکوٰۃ لینے والا
مشکل سے ملتا تھا۔
\end{hal}

%---------------------------------------------------------------------
\begin{sawal}{سوال \textenglish{14} \ \textnormal{\small(بہار \textenglish{2012}، سوال \textenglish{5} اور بہار \textenglish{2013}، سوال \textenglish{5})}}
بنو عباس کا تعارف تحریر کریں، ہارون الرشید کے کارنامے بیان کریں، نیز عہدِ بنو
عباس کی علمی و ادبی سرگرمیوں کا جائزہ لیں۔
\end{sawal}

\begin{hal}
\textbf{(الف) تعارف۔} خلافتِ عباسیہ: \textenglish{132--656}ھ
(\textenglish{750--1258}ء) --- تقریباً پانچ سو سال، اسلامی تاریخ کی سب سے
طویل خلافت۔ نسب حضرت عباسؓ، چچائے رسول \saw سے ملتا ہے، اسی سے نام پڑا۔

\smallskip
\textbf{عروج کے اسباب:} امویوں سے موالی اور اہلِ عجم کی ناراضی \ \textbullet\ \
اہلِ بیت کے نام پر چلائی گئی تحریک \ \textbullet\ \ خراسان میں ابو مسلم خراسانی
کی تنظیم \ \textbullet\ \ اموی خاندان کا اندرونی انتشار۔

\smallskip
\textbf{اہم خلفاء:} ابو العباس عبد اللہ ''السفاح'' (پہلا خلیفہ) \ \textbullet\ \
ابو جعفر منصور --- جس نے \textenglish{145}ھ (\textenglish{762}ء) میں
\textbf{بغداد} کی بنیاد رکھی اور خلافت کو حقیقی استحکام دیا \ \textbullet\ \
ہارون الرشید \ \textbullet\ \ مامون الرشید \ \textbullet\ \ اور آخری خلیفہ
المستعصم باللہ۔

\medskip
\textbf{(ب) ہارون الرشید کے کارنامے۔} دورِ خلافت: \textenglish{170--193}ھ
(\textenglish{786--809}ء) --- یہی عباسی خلافت کا \textbf{عہدِ زریں} ہے۔

\begin{itemize}\setlength{\itemsep}{1pt}
  \item \textbf{بیت الحکمہ} کی بنیاد --- کتب خانہ، ترجمہ گھر اور رصد گاہ؛ کمال
        تک ان کے بیٹے مامون الرشید نے پہنچایا۔
  \item \textbf{عدل} --- امام ابو یوسفؒ کو قاضی القضاۃ مقرر کیا؛ انہی کی
        فرمائش پر ''کتاب الخراج'' لکھی گئی، اسلامی مالیات کی بنیادی کتاب۔
  \item \textbf{سفارت کاری} --- یورپ کے شہنشاہ شارلمین سے تحائف و تعلقات۔
  \item \textbf{رفاہِ عامہ} --- ہسپتال، مساجد، مدارس؛ ملکہ زبیدہ کا مکہ کے لیے
        پانی کا انتظام، آج تک \emph{نہرِ زبیدہ} کے نام سے مشہور۔
  \item \textbf{کمزور پہلو} --- خاندانِ برامکہ کا اچانک خاتمہ۔
\end{itemize}

\medskip
\textbf{(ج) علمی و ادبی سرگرمیاں۔} یہ سوال بہار \textenglish{2013} میں مستقل
سوال کے طور پر آیا۔ جواب کو \emph{فنون کے عنوانات} میں تقسیم کیجیے --- یہی سب
سے زیادہ نمبر دلاتا ہے۔

\begin{itemize}\setlength{\itemsep}{1pt}
  \item \textbf{تحریکِ ترجمہ اور کاغذ سازی} --- یونانی و فارسی کتب کا عربی
        ترجمہ (حنین بن اسحاق)؛ سمرقند سے بغداد پہنچی کاغذ سازی نے نقل سستی
        بنائی۔
  \item \textbf{فقہ و حدیث} --- چاروں امام (ابو حنیفہؒ، مالکؒ، شافعیؒ، احمد بن
        حنبلؒ) اور \textbf{صحاح ستہ کا پورا مجموعہ} اسی عہد میں مرتب ہوا۔
  \item \textbf{ریاضی، سائنس و ادب} --- الخوارزمی (\textenglish{algorithm}
        کا لفظ انہی کے نام سے)، جابر بن حیان، ابن سینا، البیرونی؛ المتنبی،
        الفارابی۔
\end{itemize}

\smallskip
\textbf{خلاصہ۔} اموی عہد نے سلطنت دی، عباسی عہد نے \emph{تہذیب} دی۔
\end{hal}

%---------------------------------------------------------------------
\begin{sawal}{سوال \textenglish{15} \ \textnormal{\small(بہار \textenglish{2012}، سوال \textenglish{6} اور خزاں \textenglish{2021}، سوال \textenglish{8})}}
ہلاکو خان کے حملۂ بغداد پر نوٹ لکھیں، نیز ''اندلس میں اسلام'' اور ''مرابطون''
پر مختصر نوٹ تحریر کریں۔
\end{sawal}

\begin{hal}
\textbf{(الف) سقوطِ بغداد --- \textenglish{656}ھ / \textenglish{1258}ء۔}
حملہ آور: \textbf{ہلاکو خان}، چنگیز خان کا پوتا۔ آخری عباسی خلیفہ:
\textbf{المستعصم باللہ}۔

\smallskip
\textbf{اسباب:} خلافت کی عسکری و سیاسی کمزوری \ \textbullet\ \ صوبوں کی
خودمختاری اور مرکز کا بے اثر ہو جانا \ \textbullet\ \ اندرونی فرقہ وارانہ
کشیدگی \ \textbullet\ \ خلیفہ کی بروقت فیصلہ نہ کر سکنے کی عادت \ \textbullet\ \
اور روایتی بیان کے مطابق وزیر ابن العلقمی کی غداری --- اگرچہ جدید مؤرخین اس
الزام کو مشکوک سمجھتے ہیں، اس لیے اسے حتمی نہ لکھیے۔

\smallskip
\textbf{واقعات و نتائج:} چالیس روزہ محاصرے کے بعد شہر فتح ہوا؛ بے پناہ قتلِ عام
ہوا؛ کتب خانے اور \textbf{بیت الحکمہ} تباہ کر دیے گئے --- روایت ہے کہ دجلہ کا
پانی کتابوں کی سیاہی سے سیاہ ہو گیا۔ پانچ سو سالہ عباسی خلافت خاتمے کو پہنچی
اور علمی مرکز ہمیشہ کے لیے ختم ہو گیا۔

\smallskip
\textbf{اہم بات:} منگولوں کا سیلاب دو سال بعد \textenglish{658}ھ
(\textenglish{1260}ء) میں \textbf{عین جالوت} کے مقام پر مصر کے مملوکوں (سلطان
قطز اور بیبرس) کے ہاتھوں رکا۔ یہ نکتہ لکھنا جواب کو مکمل کرتا ہے۔

\medskip
\textbf{(ب) اندلس میں اسلام۔}
\begin{itemize}\setlength{\itemsep}{1pt}
  \item \textbf{آغاز} --- \textenglish{92}ھ (\textenglish{711}ء): طارق بن زیاد
        نے وادیِ لکہ کی جنگ میں راڈرک کو شکست دی۔
  \item \textbf{اموی امارت} --- عبد الرحمن الداخل نے \textenglish{138}ھ
        (\textenglish{756}ء) میں قرطبہ میں آزاد امارت قائم کی۔
  \item \textbf{خلافت کا اعلان} --- عبد الرحمن الناصر نے \textenglish{316}ھ
        (\textenglish{929}ء) میں قرطبہ میں خلافت کا اعلان کیا --- اندلس کا
        عروج۔
  \item \textbf{تہذیبی خدمات} --- جامع مسجد قرطبہ، الحمرا (غرناطہ)؛ ابن رشد،
        ابن حزم، الزہراوی۔ یورپ کی نشاۃ الثانیہ کا بڑا حصہ اسی علم سے آیا۔
  \item \textbf{انجام} --- \textenglish{897}ھ (\textenglish{1492}ء) میں غرناطہ
        کے سقوط کے ساتھ آٹھ سو سالہ اسلامی اندلس کا خاتمہ۔
\end{itemize}

\medskip
\textbf{(ج) مرابطون۔} \textenglish{448--541}ھ (\textenglish{1056--1147}ء) ---
شمالی افریقہ و اندلس کی ایک بربر (صنہاجہ) سلطنت۔ ''رباط'' (دینی و عسکری
تربیتی مراکز) سے وابستہ لوگ ''مرابطون'' کہلائے۔ بانی: عبد اللہ بن یاسین
(دینی رہنما) اور یوسف بن تاشفین، جنہوں نے \textenglish{454}ھ میں
\textbf{مراکش} کی بنیاد رکھی۔ \textbf{سب سے اہم کارنامہ --- جنگِ زلاقہ}
(\textenglish{479}ھ / \textenglish{1086}ء): اندلس کی کمزور مسلم ریاستوں کی
درخواست پر یوسف بن تاشفین نے الفانسو ششم کو فیصلہ کن شکست دی، جس نے مسلم
اقتدار کو مزید ڈیڑھ سو سال دے دیے۔ \textenglish{541}ھ میں موحدون نے ان کی جگہ
لے لی۔
\end{hal}
